All three clauses are Fact \claimref{rungs.imported-totality-minimality-completion} applied to $(\R^{1,3}, \eta)$, which is nondegenerate of mixed signature: the no-go for element-level inverses, the minimality and uniqueness of the one-plane completion, and the totality of the reciprocal on the quadric. The rung is determined by the quadratic space alone, whose null cone contains every $p^{i}$, and the reciprocal's totality is a statement about all null vectors of that cone simultaneously. A second conformal spinor changes the configuration, not the quadratic space. The corollary adds nothing to Fact \claimref{rungs.imported-totality-minimality-completion} beyond the observation that its hypotheses concern the quadratic space and not the number of null vectors named. The script confirms, in the blade engine, that a null $p$ stays null in $\Cl(2,4)$ and in $\Cl(3,5)$ (the rung is an isometric extension and inverts nothing as an element) and that the rungs supply one further anticommuting null pair each, orthogonal to $p$.
